% !TEX root = ../journal_0421052099.tex
%
% Condensed from thesis_book/chapters/conclusion.tex.
\section{Conclusion}\label{sec:conclusion}

This paper asked what an agent gains by navigating a knowledge graph rather than
retrieving from it. It also asked what a claim-level check against the traversed
subgraph adds on top. AGR leads every accuracy column on WebQSP and CWQ, at half
the language-model calls of the agentic baseline. Its margin over that baseline
is affordability under a fixed budget rather than better resolution. On CWQ it
is the only system whose accuracy rises with hop count. Structural hallucination
is eliminated by navigation itself. The verification layer contributes
auditability rather than accuracy, and decomposition helps or hurts depending on
the structure of the question.

The failure census earns three small repairs, each backed by a repeated
mechanism. The claim decomposer should be told which argument of the drafted
sentence is the answer. The evaluator's decisions should be logged beside the
scores of the relations it abandons. And the drafter should commit to a value
that the evaluator resolved and the verifier grounded. One logging change has
the highest value for its cost. Persisting each accepted claim with the triples
that matched it would turn wrongful acceptance into a rate with a denominator.
It would also make the output contract checkable after the run. Three further
additions are specified by the failure categories. A decomposition step should
be gated on whether its rewrite keeps the question's interrogative and its
constraints. A set-intersection operator would serve the $47$ composite-claim
failures, and ordinal and literal-value support would serve the knowledge-graph
gaps. A semantic check that accepts a bounded relation chain would recover the
claims wrongly rejected for lacking a single edge. Finally, re-running the
agentic baseline at AGR's candidate widths, and the static baseline at a wider
and ordered radius, would lift the two baseline bounds.

Four directions extend the work rather than repairing it. A knowledge graph
constructed by a language model, used as a controlled comparison, would separate
what the framework contributes from what a curated environment contributes.
Rollout-based value estimation would supply what the scorer lacks, since it
scores each candidate once and never revises that score. Temporal and
multi-source environments would address the questions with a time qualifier that
recur through the census. And high-stakes domains are where auditability matters
more than accuracy.

The result we would most want carried forward is not the accuracy margin. It is
the shape of the gap that remains. A system that traverses a real graph and
copies its answers out of real triples can be made to assert nothing ungrounded.
What remains is an answer that is grounded, connected, and true, and is still
not the answer to the question that was asked. Fact verification against a
knowledge graph turns out to be the easier half of the problem. \emph{Relevance}
verification is the half still open.

\endinput
